\originalchapter{散射截面、衰变率与 QED 基本过程}
\label{ch:qed-processes}

Feynman 图给出复振幅;实验计数的却是在指定时间、能量和角度范围内发生的事件数.把二者连接起来,需要处理外态归一化、能量动量守恒、可用末态数和入射通量.我们先完成这一转换,再把 Dirac 自旋求和和光子规范检查用到两个具体过程:电子正电子湮灭和 Compton 散射.后者还给出检验量子理论的经典低能极限.

\originalsection{振幅与 Lorentz 不变相空间}

沿用协变外态 $\ket p_{\rm cov}=\sqrt{2E_{\mathbf p}}\ket p$,定义连通散射矩阵元
\begin{equation}
 {}_{\rm cov}\bra f(S-I)\ket i_{\rm cov}
 =\ii(2\pi)^4\delta^4(P_f-P_i)\mathcal M_{fi}.
 \label{eq:invariant-amplitude}
\end{equation}
$\mathcal M$ 不含总体动量守恒 delta.下文 $\overline{|\mathcal M|^2}$ 表示对实际未偏振入射态取平均、对未测量末态自旋或偏振求和;若实验选择偏振,则改用相应的态密度矩阵,不能继续平均.

为了理解 delta 的平方,先在体积 $V$、观测时间 $T$ 内归一化平面波.由 Fourier 积分,$(2\pi)\delta(0)=T$,$(2\pi)^3\delta^3(0)=V$,所以
\begin{equation}
 [(2\pi)^4\delta^4(P)]^2=VT\,(2\pi)^4\delta^4(P).
\end{equation}
这应理解为有限时间和有限体積极限中的分布规则,而非点值乘法.每条外腿的盒归一化因子为 $(2EV)^{-1/2}$;每个末态动量格点求和变成 $V\int\dd^3p/(2\pi)^3$.把振幅平方、末态求和,再除以 $T$,可见所有末态体积因子消去,留下两入射粒子的一个 $1/V$.

\begin{definition}[相空间]
对总入射四动量 $P$ 和 $n$ 个质量为 $m_j$ 的末态粒子,定义
\begin{equation}
 \dd\Phi_n(P)=(2\pi)^4\delta^4\!\left(P-\sum_{j=1}^np_j\right)
 \prod_{j=1}^n\frac{\dd^3p_j}{(2\pi)^3\,2E_j},
 \qquad E_j=\sqrt{\mathbf p_j^2+m_j^2}.
 \label{eq:phase-space}
\end{equation}
\end{definition}

$\dd^3p/(2E)$ 的不变性可由 $\dd^4p\,\delta(p^2-m^2)\theta(p^0)$ 对 $p^0$ 积分看出.故相空间整体不依赖参考系.若末态有 $n_a$ 个同种不可区分粒子,且积分对它们的标签全空间独立计数,则乘 $1/\prod_a n_a!$;否则一个物理末态被数了多次.这个因子与振幅内部的费米反对称化是不同步骤.

两入射粒子的盒中跃迁率为
\begin{equation}
 \dd W=\frac{1}{4E_1E_2V}\overline{|\mathcal M|^2}\dd\Phi_n.
\end{equation}
定义通量速度
\begin{equation}
 v_{\rm flux}=\frac{\sqrt{(p_1\cdot p_2)^2-m_1^2m_2^2}}{E_1E_2}.
\end{equation}
在一粒子静止系它是另一粒子的速度,在共线迎头碰撞中是两速率之和;一般非共线时不要把它误称为普通速度差的模.单位体积中的两粒子通量为 $v_{\rm flux}/V$,由率除以通量得到
\begin{equation}
 \dd\sigma=
 \frac{\overline{|\mathcal M|^2}}{4\sqrt{(p_1\cdot p_2)^2-m_1^2m_2^2}}
 \dd\Phi_n.
 \label{eq:cross-section-general}
\end{equation}
同种末态因子按上述规则另乘.四维截面的质量量纲为 $-2$,这个量纲是检查漏掉能量因子的方便方法.

\originalsection{二体相空间与衰变率}

在质心系 $P=(\sqrt s,\mathbf0)$,先用空间 delta 消去 $\mathbf k_2$,再写 $\dd^3k_1=k^2\dd k\dd\Omega$,得到
\begin{equation}
 \dd\Phi_2=\frac{1}{16\pi^2}\frac{k^2\dd k\dd\Omega}{E_3E_4}
 \delta(\sqrt s-E_3-E_4).
\end{equation}
零点处 $\frac{\dd(E_3+E_4)}{\dd k}=k\sqrt s/(E_3E_4)$,因此
\begin{equation}
 \dd\Phi_2=\frac{k_*}{16\pi^2\sqrt s}\dd\Omega,
 \quad k_*=\frac{\sqrt{\lambda_K(s,m_3^2,m_4^2)}}{2\sqrt s},
 \quad \lambda_K(a,b,c)=a^2+b^2+c^2-2ab-2ac-2bc.
 \label{eq:two-body-phase}
\end{equation}
$\lambda_K$ 是 Källén 多项式,不是标量场的耦合 $\lambda$.质心系的入射通量根号为 $p_*\sqrt s$,于是
\begin{equation}
 \frac{\dd\sigma}{\dd\Omega}
 =\frac{1}{64\pi^2s}\frac{k_*}{p_*}\overline{|\mathcal M|^2}.
 \label{eq:two-body-cross}
\end{equation}
末态阈值出现于 $\sqrt s=m_3+m_4$,其附近的相空间压缩由 $k_*$ 明确给出.

一个质量为 $M$ 的粒子在静止系衰变时,外态归一化同样导出
\begin{equation}
 \dd\Gamma=\frac{1}{2M}\overline{|\mathcal M|^2}\dd\Phi_n.
\end{equation}
这里没有第二个入射粒子,所以没有两粒子通量分母.在其他系中每单位坐标时间的衰变率为 $(M/E)\Gamma$,与时间膨胀一致.

\begin{example}
实标量 $X$ 与实标量 $\phi$ 通过 $\Lag_{\rm int}=-gX\phi^2/2$ 相互作用,质量分别为 $M,m$,且 $M>2m$.求最低阶 $X\to\phi\phi$ 的总衰变率.
\end{example}
\begin{solution}
两条相同 $\phi$ 腿抵消顶点中的 $2!$,故 $\ii\mathcal M=-\ii g$,$|\mathcal M|^2=g^2$.用 $k_*=(M/2)\sqrt{1-4m^2/M^2}$ 和 $\int\dd\Phi_2=k_*/(4\pi M)$,并乘相同末态 $1/2!$,有
\[
 \Gamma=\frac1{2M}\frac12 g^2\frac{k_*}{4\pi M}
 =\frac{g^2}{32\pi M}\sqrt{1-\frac{4m^2}{M^2}}.
\]
顶点中的 $1/2$ 与相空间中的 $1/2!$ 各有来源,不能因为前者已经使用就删掉后者.
\end{solution}

\originalsection{电子正电子湮灭的振幅与自旋求和}

考虑 $e^-(p_1)e^+(p_2)\to\mu^-(k_1)\mu^+(k_2)$.两个 Dirac 场各自与同一光子耦合,最低阶只有一张 $s$ 道虚光子图.为让计算主线清楚,忽略电子质量,保留缪子质量 $M=m_\mu$;近电子阈值时不能使用此近似.定义 Mandelstam 变量
\begin{equation}
 s=(p_1+p_2)^2,\qquad t=(p_1-k_1)^2,
 \qquad u=(p_1-k_2)^2,
 \qquad s+t+u=2M^2.
\end{equation}
由两个 $-\ii e\gamma^\mu$ 顶点和 Feynman 规范光子传播子,得到
\begin{equation}
 \mathcal M=\frac{e^2}{s}
 [\bar v(p_2)\gamma^\mu u(p_1)]
 [\bar u(k_1)\gamma_\mu v(k_2)].
 \label{eq:ee-mumu-amplitude}
\end{equation}
总体相位不影响截面.入射电子流与 $q=p_1+p_2$ 收缩时为零,缪子流也一样;这里反粒子方程是 $(\slashed p+m)v=0$,正好使质量项消去.因此虚光子纵向规范项不贡献.

未偏振入射共有 $2\times2$ 种自旋,平均因子为 $1/4$.用自旋完备关系,
\begin{align}
 \overline{|\mathcal M|^2}
 &=\frac{e^4}{4s^2}L^{\mu\nu}H_{\mu\nu},\\
 L^{\mu\nu}&=\tr[\slashed p_2\gamma^\mu\slashed p_1\gamma^\nu]
 =4[p_2^\mu p_1^\nu+p_2^\nu p_1^\mu-(p_1\cdot p_2)\eta^{\mu\nu}],\\
 H_{\mu\nu}&=\tr[(\slashed k_1+M)\gamma_\mu(\slashed k_2-M)\gamma_\nu]
 =4[k_{1\mu}k_{2\nu}+k_{1\nu}k_{2\mu}
 -(k_1\cdot k_2+M^2)\eta_{\mu\nu}].
\end{align}
$H$ 中质量平方项是负号;把反粒子自旋求和错写成 $\slashed k_2+M$ 会在此改变阈值行为.将两组对称张量相乘,含 $\eta_{\mu\nu}$ 的项相互抵消一部分,留下
\begin{equation}
 L^{\mu\nu}H_{\mu\nu}=32[(p_2\cdot k_1)(p_1\cdot k_2)
 +(p_2\cdot k_2)(p_1\cdot k_1)+M^2(p_1\cdot p_2)].
\end{equation}
再代入 $p_1\cdot p_2=s/2$、$p_1\cdot k_1=(M^2-t)/2$、$p_1\cdot k_2=(M^2-u)/2$,其他两个点积由守恒相同,得到
\begin{equation}
 \overline{|\mathcal M|^2}
 =\frac{2e^4}{s^2}[(t-M^2)^2+(u-M^2)^2+2M^2s].
 \label{eq:ee-mumu-square}
\end{equation}

\originalsection{角分布、阈值与高能极限}

在质心系,取入射电子沿 $z$ 轴,末态缪子与它夹角为 $\theta$.定义 $\beta=\sqrt{1-4M^2/s}$,于是 $p_*=\sqrt s/2$、$k_*=\beta\sqrt s/2$,且
\begin{equation}
 t-M^2=-\frac s2(1-\beta\cos\theta),
 \qquad u-M^2=-\frac s2(1+\beta\cos\theta).
\end{equation}
代入式 \eqref{eq:ee-mumu-square},振幅平方变为
\begin{equation}
 \overline{|\mathcal M|^2}=e^4[1+\cos^2\theta+(1-\beta^2)\sin^2\theta].
\end{equation}
因此
\begin{equation}
 \frac{\dd\sigma}{\dd\Omega}
 =\frac{\alpha^2\beta}{4s}
 [1+\cos^2\theta+(1-\beta^2)\sin^2\theta],
 \qquad\alpha=\frac{e^2}{4\pi}.
 \label{eq:mumu-angular}
\end{equation}
用 $\int\dd\Omega=4\pi$、$\int\cos^2\theta\dd\Omega=4\pi/3$、$\int\sin^2\theta\dd\Omega=8\pi/3$,得到
\begin{equation}
 \sigma=\frac{4\pi\alpha^2}{3s}\,
 \beta\left(1+\frac{2M^2}{s}\right).
\end{equation}
高能 $s\gg M^2$ 时为 $4\pi\alpha^2/(3s)$,角分布为 $1+\cos^2\theta$.阈值处振幅平方趋向常数,截面仍因可用相空间 $\beta\to0$ 而消失;角分布趋向各向同性.极靠近阈值时末态电磁相互作用的高阶效应可能需要重新求和,上式只是树级结果.

\originalsection{Compton 双图与规范检验}

考虑 $e^-(p)+\gamma(k)\to e^-(p')+\gamma(k')$,其中 $p^2=p'^2=m^2$、$k^2=k'^2=0$,$p+k=p'+k'$.沿电子线,入射和出射光子有两种顶点次序,所以必须相加两张图.记 $\epsilon,\epsilon'^*$ 为对应偏振,振幅为
\begin{equation}
 \mathcal M=-e^2\bar u(p')\left[
 \slashed\epsilon'^*\frac{\slashed p+\slashed k+m}{2p\cdot k}\slashed\epsilon
 +\slashed\epsilon\frac{\slashed p-\slashed k'+m}{-2p\cdot k'}\slashed\epsilon'^*
 \right]u(p).
 \label{eq:compton-amplitude}
\end{equation}
分母由内部电子 $q_s=p+k$、$q_u=p-k'$ 的 $q^2-m^2$ 得到.两项之间是加号;它们是同一条电子线的两个顶点次序,没有交换两个相同外电子.

\begin{center}
\begin{tikzpicture}[x=1cm,y=1cm,>=Stealth,
 photon/.style={decorate,decoration={snake,amplitude=1.2pt,segment length=5pt}},
 fermion/.style={->,thick}]
 \draw[fermion] (-2,0)--(-.6,0); \draw[fermion] (-.6,0)--(.6,0); \draw[fermion] (.6,0)--(2,0);
 \draw[photon] (-1.4,1)--(-.6,0);\draw[photon] (.6,0)--(1.4,1);
 \node at (-1.5,-.35){$p$};\node at (1.5,-.35){$p'$};\node at (-1.5,1.25){$k$};\node at (1.5,1.25){$k'$};
 \node at (0,-.4){$p+k$};
 \begin{scope}[xshift=5.6cm]
 \draw[fermion] (-2,0)--(-.6,0); \draw[fermion] (-.6,0)--(.6,0); \draw[fermion] (.6,0)--(2,0);
 \draw[photon] (-.6,0)--(-1.4,1);\draw[photon] (1.4,1)--(.6,0);
 \node at (-1.5,-.35){$p$};\node at (1.5,-.35){$p'$};\node at (-1.5,1.25){$k'$};\node at (1.5,1.25){$k$};
 \node at (0,-.4){$p-k'$};
 \end{scope}
\end{tikzpicture}
\end{center}

检验入射光子 $\epsilon\mapsto k$.令 $D(q)=\slashed q-m$,则 $D(q)^{-1}=(\slashed q+m)/(q^2-m^2)$.第一图的 $\slashed k=D(p+k)-D(p)$;由于 $D(p)u(p)=0$,第一项退化为 $\bar u(p')\slashed\epsilon'^*u(p)$.第二图利用 $p-k'=p'-k$,故 $\slashed k=D(p')-D(p-k')$;左边 $\bar u(p')D(p')=0$,第二项退化为前者的负值.二者相消.对出射偏振 $\epsilon'^*\mapsto k'$ 用同一恒等式也得到零.单张图各自不为零,必须保留双图才满足规范不变性.

\originalsection{Compton 自旋求和的主计算链}

为了取得未偏振截面,利用上一节的规范检查,用 $-\eta_{\mu\nu}$ 替代两处物理偏振和.设 $A=p\cdot k>0$、$B=p\cdot k'>0$,并令
\begin{equation}
 X^{\nu\mu}=\gamma^\nu\frac{\slashed q_s+m}{2A}\gamma^\mu
 -\gamma^\mu\frac{\slashed q_u+m}{2B}\gamma^\nu.
\end{equation}
$\bar X^{\nu\mu}=\gamma^0(X^{\nu\mu})^\dagger\gamma^0$ 表示矩阵乘积反向后的 Dirac 共轭;在实动量下仅反转 gamma 次序.于是
\begin{equation}
 \overline{|\mathcal M|^2}=\frac{e^4}{4}
 \tr[(\slashed p'+m)X^{\nu\mu}(\slashed p+m)\bar X_{\nu\mu}].
 \label{eq:compton-trace}
\end{equation}
$1/4$ 是入射电子二自旋与光子二偏振的平均.为避免把八矩阵迹藏在一句结论中,下面把它拆成两项平方和一项干涉.由守恒,
\begin{equation}
 q_s^2=m^2+2A,\quad q_u^2=m^2-2B,
 \quad p\cdot p'=m^2+A-B,\quad k\cdot k'=A-B,
 \quad p'\cdot q_s=m^2+A.
\end{equation}
平方 $s$ 项先收缩最内层 $\gamma^\mu(\slashed p+m)\gamma_\mu=-2\slashed p+4m$.再用
$\slashed q\slashed p\slashed q=2(q\cdot p)\slashed q-q^2\slashed p$,有
\begin{equation}
 (\slashed q_s+m)(-2\slashed p+4m)(\slashed q_s+m)
 =4[(m^2-A)\slashed q_s+A\slashed p+m(m^2+A)].
\end{equation}
外层 gamma 收缩使矢量部分乘 $-2$、标量部分乘 $4$,最终迹为
\begin{equation}
 T_{ss}=32[m^4+m^2A+AB].
\end{equation}
因此 $s$ 平方项是 $e^4T_{ss}/(16A^2)=2e^4[B/A+m^2/A+m^4/A^2]$.交换入、出光子的交叉替换 $A\mapsto-B$、$B\mapsto-A$,同样得到 $u$ 平方项 $2e^4[A/B-m^2/B+m^4/B^2]$.

干涉的一个迹为
\begin{equation}
 T_{su}=\tr[(\slashed p'+m)\gamma^\nu(\slashed q_s+m)\gamma^\mu
 (\slashed p+m)\gamma_\nu(\slashed q_u+m)\gamma_\mu].
\end{equation}
令 $a=p'$、$b=q_s$、$c=p$、$d=q_u$,按留下的质量因子数把迹写成 $T_{su}=T^{(0)}+T^{(2)}+T^{(4)}$;奇数 gamma 迹为零.无质量项先用
$\gamma^\mu\slashed c\gamma_\nu\slashed d\gamma_\mu=-2\slashed d\gamma_\nu\slashed c$,再用 $\gamma^\nu\slashed b\slashed d\gamma_\nu=4b\cdot d$,得
\begin{align}
 T^{(0)}&=-32(b\cdot d)(a\cdot c)=-32m^2(m^2+A-B),\\
 T^{(2)}&=16m^2(a\cdot b+a\cdot c+a\cdot d+b\cdot c+b\cdot d+c\cdot d),\\
 T^{(4)}&=-32m^4.
\end{align}
第二行来自选择两处 $m$ 的六个项;每项按相同收缩公式化为 $16m^2$ 乘剩下的两个动量点积.六个点积依次为 $m^2+A$、$m^2+A-B$、$m^2-B$、$m^2+A$、$m^2$、$m^2-B$,所以 $T^{(2)}=16m^2[6m^2+3(A-B)]$.三组相加得到
\begin{equation}
 T_{su}=16m^2(A-B+2m^2).
\end{equation}
它的共轭相同,两个干涉项的分母符号都为负,故合计是 $-e^4T_{su}/(8AB)$.相加三部分,$m^4/A^2,m^4/B^2$ 与干涉中的 $-2m^4/(AB)$ 组合成平方,得到
\begin{equation}
 \overline{|\mathcal M|^2}=2e^4\left[
 \frac AB+\frac BA+2m^2\left(\frac1A-\frac1B\right)
 +m^4\left(\frac1A-\frac1B\right)^2\right].
 \label{eq:compton-invariant-square}
\end{equation}
各单图用协变偏振求和得到的中间值本身不是独立物理概率;例如低能处的大项只有相加后才抵消.这也解释了为何干涉项不能忽略.

\originalsection{实验室截面与 Thomson 极限}

取入射电子静止,$p=(m,\mathbf0)$,光子能量为 $\omega,\omega'$,两光子方向夹角为 $\theta$.由末态电子壳条件
$(p+k-k')^2=m^2$,得
\begin{equation}
 m(\omega-\omega')=\omega\omega'(1-\cos\theta),
 \qquad \omega'=\frac{\omega}{1+(\omega/m)(1-\cos\theta)}.
 \label{eq:compton-energy}
\end{equation}
相空间中的能量 delta 为 $\delta(m+\omega-\omega'-E')$,$E'=\sqrt{m^2+|\mathbf k-\mathbf k'|^2}$.在零点处其导数绝对值为
$(E'+\omega'-\omega\cos\theta)/E'=[m+\omega(1-\cos\theta)]/E'=m\omega/(E'\omega')$.因此对 $\omega'$ 积分后
\begin{equation}
 \dd\Phi_2=\frac1{16\pi^2}\frac{\omega'^2}{m\omega}\dd\Omega,
 \qquad
 \frac{\dd\sigma}{\dd\Omega}
 =\frac{\overline{|\mathcal M|^2}}{64\pi^2m^2}
 \left(\frac{\omega'}\omega\right)^2.
\end{equation}
这里入射通量为 $4m\omega$.又有 $A=m\omega$、$B=m\omega'$,以及 $m^2(1/A-1/B)=-(1-\cos\theta)$.代入式 \eqref{eq:compton-invariant-square},得到 Klein--Nishina 公式
\begin{equation}
 \frac{\dd\sigma}{\dd\Omega}
 =\frac{\alpha^2}{2m^2}\left(\frac{\omega'}\omega\right)^2
 \left(\frac\omega{\omega'}+\frac{\omega'}\omega-\sin^2\theta\right).
\end{equation}
当 $\omega\ll m$ 时,电子反冲能量很小,$\omega'/\omega\to1$,故
\begin{equation}
 \frac{\dd\sigma_T}{\dd\Omega}
 =\frac{\alpha^2}{2m^2}(1+\cos^2\theta),
 \qquad\sigma_T=\frac{8\pi\alpha^2}{3m^2}.
\end{equation}
这与经典 Thomson 散射一致.其经典来源也可直接理解:非相对论电荷在入射电场作用下满足 $m\ddot{\mathbf x}=e\mathbf E$,加速辐射的方向因子是横向投影.对两种入射偏振平均、两种出射偏振求和,$(1/2)\sum|\boldsymbol\epsilon'^*\cdot\boldsymbol\epsilon|^2=(1+\cos^2\theta)/2$;在有理化单位中辐射长度尺度 $e^2/(4\pi m)=\alpha/m$,于是恢复同一角分布与总截面.量子公式不仅给出这个低能结果,还给出经典无反冲近似无法包含的能量转移式 \eqref{eq:compton-energy}.
